%% appendix/_template_appendix.tex
%% Plantilla para apéndices — Control Automático UTN IEE
%%
%% INSTRUCCIONES:
%%   1. Copiar a appendix/apXX_nombre.tex
%%   2. En main.tex, agregar DENTRO del bloque \appendix:
%%      \include{appendix/apXX_nombre}
%%   3. Los apéndices usan numeración A, B, C... automáticamente.

\chapter{[Título del Apéndice]}
\label{ap:[etiqueta]}

[Breve descripción de qué contiene este apéndice y a qué capítulos está asociado.]

%======================================================================
\section{[Sección del apéndice]}
\label{sec:ap[etiqueta]-1}
%======================================================================

[Contenido: demostraciones largas, tablas de referencia, derivaciones
alternativas, material complementario que sería demasiado extenso en el
cuerpo del capítulo.]

%% Usos típicos de apéndices en este apunte:
%%
%% Ap. A: Tabla de transformadas de Laplace (cap. 2)
%% Ap. B: Derivación completa de ecuaciones de estado de máquinas eléctricas (cap. 14)
%% Ap. C: Tablas de respuesta en frecuencia de sistemas estándar (cap. 9-10)
%% Ap. D: Código MATLAB completo con comentarios extendidos
%% Ap. E: Demostraciones de criterios de estabilidad (Routh, Nyquist)
%% Ap. F: Parámetros de motores y generadores usados en los ejemplos
